\par\Needspace{6\baselineskip}\medskip
\begin{center}\hypertarget{native-172-article125}{}\pdfbookmark[2]{125.}{article-125}125.\end{center}
Durch eine geringe Abwandlung des vorhergehenden Verfahrens gewinnen wir ein \emph{zweites}. Wie dort gehen wir von den Entfernungen bei zwei Beobachtungen aus und bestimmen alle Elemente. Für die dritte Beobachtung führen wir die Rechnung aus diesen Elementen nur bis zum heliozentrischen Ort in der Bahn und lassen den geozentrischen Ort unberechnet. Andererseits leiten wir denselben heliozentrischen Ort aus dem beobachteten geozentrischen Ort und der Lage der Bahnebene her, indem wir die in den Artikeln 74 und 75 behandelte Aufgabe lösen. Diese beiden Bestimmungen weichen voneinander ab, sofern nicht zufällig die wahren Werte von $x,y$ angesetzt wurden, und liefern uns $X,Y$. Als $X$ nehmen wir die Differenz zwischen den beiden Werten der Länge in der Bahn und als $Y$ die \SourceLeaf{173}Differenz zwischen den beiden Werten des Radiusvektors oder vielmehr seines Logarithmus. Für dieses Verfahren gelten dieselben Vorbehalte wie im vorhergehenden Artikel. Ein weiterer kommt hinzu. Der heliozentrische Ort in der Bahn lässt sich aus dem geozentrischen Ort nicht herleiten, wenn der Ort der Erde mit einem der beiden Bahnknoten zusammenfällt. Dann ist dieses Verfahren also unanwendbar. Auch wenn der Ort der Erde einem der beiden Knoten sehr nahe liegt, sollte man auf dieses Verfahren verzichten. Die Annahme, dass kleinen Änderungen von $x,y$ proportionale Änderungen von $X,Y$ entsprechen, wäre aus einem ähnlichen Grund wie dem im vorhergehenden Artikel erwähnten zu ungenau. Auch hier lässt sich Abhilfe suchen, indem man den mittleren Ort mit einem der beiden äußeren vertauscht, zu dem ein von den Knoten weiter entfernter Erdort gehört, sofern die Erde nicht bei allen drei Beobachtungen zufällig in der Nähe der Knoten stand.
